\section{The terminal inputs}\label{sec:rrm-terminal-inputs}

We recall the terminal estimates used in the response proof, with the
normalizations of \cite[\S6 and Appendix~A]{TerminalPressure}.
They apply to the common trajectory fixed in the introduction.
The marked activity estimates, logarithm and complement argument
are proved in the following sections.

\subsection{Operators, norms and constants}
At the reference stop $T_* =\eta_*\Z^3/\ell_*\Z^3$,
coarse spinors live at $T_*^0=\Z^3/\ell_*\Z^3$ with weight one;
fine chamber indices have weight $\eta_*^3/48$.
The block average $Q_*$ and cluster injection $J$ are those of
\cite[\S2.1 and \S4.2]{FermionBlock}. The fine resolvent and blocked operator are
\begin{equation}\label{eq:interacting-terminal-operators}
 \begin{split}
 F_*(A)&=D_A+\mu_*+b_*JP_*(A)J^*,\qquad S_*(A)=F_*(A)^{-1},\\
 H_*(A)&=b_*S_*(A)JQ_*^T(-A),\qquad
 D_*(A)=b_*-b_*Q_*(A)J^*H_*(A).
 \end{split}
\end{equation}
Here $P_*=Q_*^T(-A)Q_*(A)$ on the cluster-averaged space,
and $H_*^\#$ includes ordinary and covariant H\"older-pair rows.
The barred rows are their charge-conjugate partners.
The parameter ranges are
\[
 b_-=1-L^{-1}\le b_*\le1,\qquad
 \mu_-=L^{-3}<\mu_*\le L^{-1},\qquad
 \gamma^FM\ge16,\qquad \ell_*>2M+2.
\]
The positive mass $\mu_*=\bar m a_*$ is distinct from the running
mass shift $m_*$ and from the possibly signed total fine mass.

The domains $\widetilde R_j$, $\mathfrak X_j$ and
$\widetilde G_j$ are the background, external and total-connection
domains of \cite[\S6]{MaxwellBlock} and \cite[\S2.3]{NormalizedContraction}.
In particular
\[
 \mathfrak X_j=\{C:|C|_\infty<e_j^{-1+2\varepsilon},\
                   |\partial C|_\infty<e_j^{-1+3\varepsilon}\},
 \qquad \widetilde R_j+\mathfrak X_j\subset\widetilde G_j.
\]
The local real-gauge definition of $\widetilde R_j$ makes it invariant
under real constants. At the stop put
\[
 \mathcal W=\{A_r+C:A_r\in\widetilde R_*\text{ real},\
                  C\in\mathfrak X_*\}.
\]
On this domain the inverse and field maps are analytic. Their
localized majorants, including short-pair H\"older rows, are bounded
by $C^F E(y,y')$, where $E(y,y')=\mathrm e^{-\gamma^Fd(y,y')}$.
For long pairs the distance is the minimum of the distances from
the two endpoint blocks. All global scalar component row sums of
$H_*^\#$ and its barred partner are at most $C^F$.

Localize a fine index in its unit block and a pair at its first
point. For an operator $O$ define
\begin{align*}
 |O|(y,y')=\sup_{\|v\|_\infty\le1}
                       \|\mathbf1_yO\mathbf1_{y'}v\|_\infty,\qquad&\\
 \|O\|_\gamma=\max\left\{
 \sup_y\sum_{y'}|O|(y,y')\mathrm e^{\gamma d(y,y')},\
 \sup_{y'}\sum_y|O|(y,y')\mathrm e^{\gamma d(y,y')}\right\}.&
\end{align*}
The spinor norm here is the maximum component modulus.
For coarse operators, $\|O\|_2\le2\|O\|_0$ and
$\|O\|_1\le4\sum_{y,y'}|O|(y,y')$, where the subscripts
$2$ and $1$ denote operator and trace norm, respectively.

For a direction $B$, let $P_B$ be the unit blocks containing
an endpoint of a bond in its support, put
$P_B^+=\{y:d(y,P_B)\le M+1\}$ and
$\epsilon_B=e_*|B|_\infty$.
The first coarse derivative has majorant
$226(C^F)^2\epsilon_B\sum_{z\in P_B}E(y,z)E(z,y')$.
The first field-leg derivative has the analogous two-edge chain
through $P_B^+$, with coefficient $C_H\epsilon_B$ and a sum
over at most two endpoint blocks. For directions separated by
$M+2$, both mixed derivatives have the sum of the two ordered
three-edge chains through $P_B,P_{B'}$, or their enlarged sets
for the legs, with coefficient $C_H^2\epsilon_B\epsilon_{B'}$.
These precise pointwise estimates are
\cite[Proposition~A.1]{TerminalPressure}; the response proof uses
their support factors as well as their row sums.

We retain their constants:
\begin{align*}
 \mathcal S&=C^F(24/\gamma^F)^3,&
 h_{\rm L}&=6\mathcal S,& C_H&=203M(C^F)^2,\\
 C_K&=226\mathcal S^2,& C_K'&=1328\mathcal S^2,&
 c_\tau&=\min\{1/8,\mu_-/48,(292\mathcal S)^{-1}\},\\
 g&=(b_-^{-1}+\mu_-^{-1})^{1/2},&h_G&=1+g,&g_1&=b_-^{-1}+4/\mu_-.
\end{align*}
Set $\theta_{\rm CT}=\mu_-/(12\mathrm e)$ and
\[
 \vartheta_0=\min\{\theta_{\rm CT}/2,\gamma^F/8,1/M\},\qquad
 \alpha_G=2b_-^{-1}
           +4\mathrm e\mu_-^{-1}(24/\theta_{\rm CT})^3.
\]
The real covariance $G=D_*(A_r)^{-1}$, for every real $A_r$,
satisfies $\|G\|_2\le g^2$ and has weighted component row and
column sums at most $\alpha_G$, with weight
$\mathrm e^{\vartheta_0d(u,u')}$.
On $\mathcal W^+=\{A+A_1:A\in\mathcal W,\ e_*|A_1|_\infty\le c_\tau\}$,
the first derivative bound holds with $C_K'$.
If $A=A_r+A_1$, $A_r$ is real, $e_*|A_1|_\infty\le\mu_*/24$
and $F_*(A)$ is invertible, then $D_*(A)$ is invertible and
$\|D_*(A)^{-1}\|_2\le g_1$.
These are \cite[(72)--(75)]{TerminalPressure}.

The activity space, independent ordinary and pair generators, and
their coefficient norm are those of \cite[\S2.2]{NormalizedContraction}.
An activity supported on an $M$-polymer $X$ has coefficients
$E_{nm}$, integrated with the fine-index weights above and pair
weights $(\eta_*^3/48)^2$, and norm
$\sum_{n,m}h_1^nh_2^m\|E_{nm}\|/(n!m!)$.
The family norm adds the factor $\mathrm e^{\kappa d_M(X)}$ and takes
the supremum over polymers and the declared connection domain.
Here $Md_M(X)$ is the infimum of the lengths of connected finite
nets meeting every closed $M$-cube of $X$. A net is a finite union
of images of closed straight segments in the continuum torus;
its length is the sum of their Euclidean lengths, with the infimum
also taken over segment representations. A point is allowed as a
segment. We use the construction constants
$\kappa_0=54\log72$, $K_0=72^{54}/36$,
$c_{\rm rb}=54K_0$ and $\kappa\ge7\kappa_0+7$.
The anchored polymer sum is bounded by $K_0$, and its reblocking
bound has constant $c_{\rm rb}$, as proved in
\cite[Appendix~B, (79)--(83)]{NormalizedContraction}.
The fixed constant $C_0\ge C^F$ bounds the photon kernel rows and
the field substitution; in particular at
$a_2=C_0^{-1}e_*^{-1/4+\varepsilon}$,
\[
 |E(H_*^\#\Psi)|_{a_2}\le\|E\|_{h_*},\qquad
 |\langle\bar\psi,\psi\rangle|_{a_2}
          \le4e_*^{-1/2+2\varepsilon}\ell_*^3,
 \quad h_*=(e_*^{-1/4},e_*^{-1/4+\varepsilon}).
\]
The scalar constant $C_d$ is the local determinant-activity bound
in \eqref{eq:interacting-determinant-input}.
All constants in this subsection are fixed before the physical charge
and are independent of the fine cutoff, twist and placement.

\subsection{The measure identity and Gaussian domination}
We use the exact terminal density of
\cite[Theorem~B]{TerminalPressure}, including its external slot.
Its original measure is $\nu^{\theta,a}$ from
\eqref{eq:external-matter-density}; write $\bar\nu$ for its
zero-source twist average and $\bar\mu=\bar\nu/\bar\nu(1)$.
The full photon-coordinate map sends an exact fine curvature
$F\in F_N=dC^1(C(n_*);\R)$ to $(A_*;\zeta_j)_{j<k_*}$.
Its Gaussian marginal at the stop is denoted $\mathsf m_*$.
Every scalar fluctuation coordinate is an orthonormal linear
functional of $F$; every coarse plaquette has row norm at most one.
The numbers of coordinates are at most
$2\ell_*^3L^{3(k_*-j)}$ at level $j$ and $3\ell_*^3$ at the stop.
The stop Gaussian is nondegenerate on the mean-zero coclosed space.

Set $\chi_*(A)=\mathbf1_{\{|dA|_\infty\le p_*\}}$ and
\[
 \mathfrak G=\{|\zeta_j(\Upsilon)|\le p_{0,j}\ (j<k_*),\
                                  |dA_*|_\infty\le p_*\}.
\]
For every bounded Borel function $\Phi$ of $A_*$,
\begin{equation}\label{eq:rrm-terminal-Borel}
 \nu^{\theta,a}(\mathbf1_{\mathfrak G}\Phi(A_*))
 =\int\Phi(A)\chi_*(A)
             Y_*^x(A+c_\theta;\mathscr C)\,\dd\mathsf m_*(A),
 \qquad \mathscr C=a^{[k_*]},\quad
 c_{\theta,\mu}=\theta_\mu/(e_*\ell_*).
\end{equation}
The density is jointly continuous on the indicated support and is
invariant under admissible real gauge changes. Its explicit fermion
formula is recalled in \cref{subsec:interacting-determinant-ratio}.
The source-free normalized measure satisfies
\begin{equation}\label{eq:rrm-zero-Gaussian}
 \bar\mu(\mathrm e^{\langle f,F\rangle})\le \mathrm e^{\|f\|^2/2}
                 \qquad(f\text{ a real two-cochain}).
\end{equation}
This is the zero-sector result
\cite[Proposition~2.1]{TerminalPressure}.
The tower and Borel identity are
\cite[Proposition~3.1, Theorem~B]{TerminalPressure}; they precede
the response argument and do not require its conclusion.
